% Generated draft: requires mathematical and source review.
\documentclass[11pt]{article}
\usepackage[margin=25mm]{geometry}
\usepackage{fontspec}
\setmainfont{Latin Modern Roman}
\usepackage{amsmath,amssymb,mathtools,unicode-math}
\setmathfont{Latin Modern Math}
\usepackage{enumitem,needspace,xurl,hyperref}
\setlength{\parindent}{0pt}
\setlength{\parskip}{5pt}
\setlength{\emergencystretch}{3em}
\newcommand{\N}{\mathbb{N}}
\newcommand{\Z}{\mathbb{Z}}
\newcommand{\Q}{\mathbb{Q}}
\newcommand{\R}{\mathbb{R}}
\newcommand{\C}{\mathbb{C}}
\newcommand{\F}{\mathbb{F}}
\newcommand{\A}{\mathbb{A}}
\newcommand{\PP}{\mathbb{P}}
\newcommand{\E}{\mathbb{E}}
\newcommand{\eps}{\varepsilon}
\newcommand{\dd}{\,\mathrm{d}}
\begin{document}
\section*{UnsolvedMath ID: 30005919}
\textbf{Source problem code:} catalogue entry
\subsection{Definitions and mathematical statement}
For \(n\in\{2,3\}\), let \(\mathbb F_n\) be the free group on \(n\) generators, and let \(\ell^2(\mathbb F_n)\) be the complex Hilbert space of square-summable functions on \(\mathbb F_n\). Its left regular representation is given by
\[
(\lambda_n(g)\xi)(h)=\xi(g^{-1}h)
\qquad(g,h\in\mathbb F_n,\ \xi\in\ell^2(\mathbb F_n)).
\]
Define the free group von Neumann algebra, called the free group factor, by
\[
L(\mathbb F_n)=\overline{\operatorname{span}_{\C}\{\lambda_n(g):g\in\mathbb F_n\}}^{\mathrm{WOT}}
\subseteq B(\ell^2(\mathbb F_n)),
\]
where \(B(H)\) denotes the bounded linear operators on \(H\), and \(\mathrm{WOT}\) denotes the weak operator topology, determined by the matrix coefficients \(T\mapsto\langle T\xi,\eta\rangle\).

\textbf{Question.} Does there exist a normal unital complex-linear \(*\)-isomorphism
\[
\Theta:L(\mathbb F_2)\longrightarrow L(\mathbb F_3)?
\]
Here an isomorphism is a bijective map preserving multiplication, adjoints, and the identity; normality means preservation of suprema of bounded increasing nets of positive operators. This expresses isomorphism as von Neumann algebras, the standard interpretation of the supplied question about free group factors.

\textbf{Unresolved source fragment.} The original statement also appends
\[
\begin{aligned}
&\phi(a_1b_1)\phi(a_2b_2)-\phi(a_1a_2b_1b_2)\\
&\qquad=\bigl(\phi(a_1b_1)-\phi(a_1)\phi(b_1)\bigr)
\bigl(\phi(a_2b_2)-\phi(a_2)\phi(b_2)\bigr).
\end{aligned}
\]
The record supplies neither definitions nor hypotheses for \(\phi,a_1,a_2,b_1,b_2\), and gives no connection between this formula and the isomorphism question. It is retained as an unexplained source fragment, without asserting its validity or imposing it as a hypothesis.
\subsection{Short English statement}
The question asks whether the complex von Neumann algebras generated by the left regular representations of the free groups on two and three generators are isomorphic. The supplied record labels the problem open; this draft does not independently verify that status. An appended formula in the original extraction lacks definitions and hypotheses.
\subsection{Sources}
\begin{itemize}
\item Local UnsolvedMath registry, \texttt{data/unsolvedmath/problems.json}: permanent ID \(30005919\), title \emph{Isomorphism of Free Group Factors}, exact source problem code \texttt{catalogue entry}; collection: \emph{Oberwolfach Reports Open Problems}. The supplied record gives proposed year 2024 and identifies no individual proposer.
\item Supplied report citation: \emph{Non-commutative Function Theory and Free Probability} (2024), Oberwolfach Reports, pp.~1227--1284, DOI 10.4171/owr/2024/22. \url{https://doi.org/10.4171/owr/2024/22}. No problem-specific page or individual attribution is supplied.
\item The registry labels the statement \texttt{corrected\_verified} and reports that it was reconstructed and checked against the official Oberwolfach report. This is supplied source-verification metadata, not an independent verification performed for this draft.
\item Supplied literature assessment, checked on \texttt{2026-08-22}: the record labels the problem open and states that the isomorphism question for free group factors of different finite ranks remains unresolved. This is supplied dated curation derived from the OWR record, not independently verified here; the record explicitly calls for independent verification of time-sensitive status and cited claims.
\item Supplied literature source labeled ``Dykema'': \url{https://arxiv.org/abs/funct-an/9211012}.
\item Supplied literature source labeled ``Popa 2024 lecture'': \url{https://www.math.ucla.edu/~popa/Talks/wFC_uclaNov2024.pdf}.
\end{itemize}
\end{document}
